\section{泛化挑战：从Narrow到General}

\subsection{Reward Model的定义难题}

朱军指出了RL泛化面临的核心挑战：

\begin{warningbox}{Reward定义的困境}
在数学定理证明和编程中，reward（奖励函数）是明确的——答案对就是对，错就是错。但在自动驾驶、艺术创作、视频生成等开放领域中，``好''和``不好''的界限模糊，每个人的感受不同。如何在这些领域定义reward model，是RL泛化的核心技术难题。
\end{warningbox}

此外，过程监督（process supervision）数据的获取也很困难——需要对思考过程的每一步进行标注，这需要专业人员和高价值数据。

\subsection{创业公司的机会}

杨植麟从创业视角分析了O1范式带来的机会：

\begin{knowledgebox}{新技术变量带来的机会}
\begin{itemize}
\item 有一定算力门槛的公司可以在RL算法层面做基础创新，甚至在基础模型上取得突破
\item 算力较少的公司可以通过后训练在特定领域做到最好
\item 有确定方向但不确定路径——对创业公司是好事
\end{itemize}
\end{knowledgebox}

\subsection{本章小结}

RL泛化的核心挑战在于开放领域的reward定义和过程监督数据获取。但技术路径已经清晰，结合强大的基座模型，泛化速度将快于上一代RL（AlphaGo时代）。


